\paragraph{\trPathSpace{}上的\trLineBundle{}与\trSpectralFlow。}%
\label{sub:path-spaces-spectral-flow}
设 $N$ 为一闭\trManifold，$\Sigma\subset N$ 为一张\trHypersurface，其上带有一个实%
\trLineBundle{} $\sigma\to\Sigma$ 以及一个取定的\trIsomorphism
\begin{equation}\label{eq:coorientation-iso}
(\sigma^{\otimes 2})^{\vee}\cong\nu_{\Sigma}.
\end{equation}
这里 $\nu_{\Sigma}$ 是\trNormalBundle，从而 $\Sigma$ 从\eqref{eq:coorientation-iso}
继承一个\trCoorientation。由这些数据可以得到：一个\trLocallyConstant{}函数
$I\colon\scr{P}\to\ZZ$，以及一个实\trLineBundle{} $\delta\to\scr{P}$，其中 $\scr{P}$
是端点落在 $\Sigma$ 之外的\trSmooth{}路径 $[0;1]\to N$ 构成的空间。用\trCohomology{}的语言来
说，这是把 \trThom{}\trIsomorphism{}与\trEvaluationMap{} $\scr{P}\times[0;1]\to N$ 结合起来完成的：
也就是说，若
\begin{equation*}
\begin{split}
S\colon H^*(\Sigma)\cong H^{*+1}(N,N\setminus\Sigma)\qquad\qquad\qquad\qquad\\
\longrightarrow H^{*+1}(\scr{P}\times[0;1],\scr{P}\times\{0;1\})\cong H^*(\scr{P}),
\end{split}
\end{equation*}
则 $I=S(1)$，类似地 $w_1(\delta)=S(w_1(\sigma))$。事实上，$I$ 是与 $[\Sigma]$ 的%
\trIntersectionNumber，因而可以在几何上解释为交点的计数。下面我们将沿同一思路解释
如何构造 $\delta$ 的一个显式代表元。
取一条与 $\Sigma$ \trTransverse{}相交的路径 $\alpha\in\scr{P}$。在有限多个\trCrossing{}
$s\in\alpha^{-1}(\Sigma)\subset(0;1)$ 的每一处，都有一个非零\trNormalVector $\alpha'(s)$，从而
由\eqref{eq:coorientation-iso} 得到 $\sigma_{\alpha(s)}$ 上的一个\trNondegenerate%
\trQuadraticForm{} $q_{\alpha}(s)$。它的符号就是与 $\Sigma$ 的局部%
\trIntersectionNumber，于是 $I(\alpha)=\sum_s\mathrm{sign}(q_{\alpha}(s))$。我们定义
\begin{equation}\label{eq:delta-alpha}
\delta_{\alpha}=\bigotimes_s\sigma_{\alpha(s)}^{\mathrm{sign}(q_{\alpha}(s))}
\end{equation}
其中 $\sigma^{-1}$ 是逆（或对偶）丛；约定\trTensorProduct{}按诸穿越时刻 $s$
的递增次序来作。要把诸 $\delta_{\alpha}$ 变成某个\trLineBundle{}的\trFiber，还需对一个%
\trGeneric{}\trHomotopy{} $\{\alpha_r\}$（$0\le r\le 1$）定义\trParallelTransport{}映射
$\delta_{\alpha_0}\to\delta_{\alpha_1}$。其中非平凡之处在于：当交
$\alpha_r\cap\Sigma$ 发生\trGeneric{}\trBifurcation{}（生或灭）时，应当如何指明所
发生的事。设在 $(\tilde{r},\tilde{s})$ 处发生一次“生”。此时 $(d/ds)^2\alpha_r$ 是
$\Sigma$ 的一个\trWellDefined{}的非零\trNormalVector，于是在
$\tilde{\sigma}=\sigma_{\alpha_{\tilde{r}}(\tilde{s})}$ 上有一个相关联的非零\trQuadraticForm
$\tilde{q}$。取充分小的 $\epsilon>0$，令 $r_{\pm}=\tilde{r}\pm\epsilon$，并考察在
时刻 $(r_{\pm},s_{\pm})$（$s_-<s_+$）出现的两个新生交点。它们将靠近\trBifurcation{}点；
若再借助一个局部\trTrivialization{}把这些点处 $\sigma$ 的\trFiber{}与 $\tilde{\sigma}$
等同起来，则 $q_{\alpha_{r_+}}(s_{\pm})$ 的符号将是 $\pm\mathrm{sign}(\tilde{q})$。
其余交点处的符号从 $\alpha_{r_-}$ 承袭而来，因此便有一个取定的\trIsomorphism
\begin{equation}\label{eq:birth-isomorphism}
\begin{split}
\delta_{\alpha_{r_+}}\cong\bigotimes_{s<\tilde{s}}
\sigma_{\alpha_{r_-}(s)}^{\mathrm{sign}(q_{\alpha_{r_-}}(s))}\otimes
\tilde{\sigma}^{-\mathrm{sign}(\tilde{q})}\otimes\tilde{\sigma}^{\mathrm{sign}(\tilde{q})}\qquad\qquad\\
\otimes\bigotimes_{s>\tilde{s}}
\sigma_{\alpha_{r_-}(s)}^{\mathrm{sign}(q_{\alpha_{r_-}}(s))}.
\end{split}
\end{equation}
取 $\tilde{v}$、$\tilde{v}^{-1}$ 分别为 $\tilde{\sigma}$ 中的一个非零向量与
$\tilde{\sigma}^{-1}$ 中相应的对偶向量，便可以定义\trParallelTransport{}映射
$\delta_{\alpha_{r_-}}\to\delta_{\alpha_{r_+}}$：把 $\tilde{v}^{-1}\otimes\tilde{v}$
或 $\tilde{v}\otimes\tilde{v}^{-1}$（视 $\tilde{q}$ 的符号而定）插入
\eqref{eq:birth-isomorphism} 之中。对“灭”的过程则使用相反的规则。要说明这确实定义了%
\trPathSpace{}上的一个\trLineBundle，还应当对“\trHomotopy{}的\trHomotopy”验证一个\trConsistency{}条件。
这很直接，我们略去不表。
原型的例子如下。令 $N=\mathrm{sym}^2(\RR^n)^{\vee}$ 为 $\RR^n$ 上\trQuadraticForm{}构成的空间，
并令 $N=\Sigma^0\cup\Sigma^1\cup\cdots$ 为其按\trCorank{}的\trStratum，其中
$\Sigma^0$ 是由\trNondegenerate{}形式组成的（开\trDense）子集。诸\trNullSpace{}构成%
\trVectorBundle{} $\sigma^d\to\Sigma^d$，并带有\trTautological{}\trIsomorphism{}
$(\mathrm{sym}^2\sigma^d)^{\vee}\cong\nu_{\Sigma^d}$。把上面的构造应用于
$\Sigma=\overline{\Sigma^1}=\Sigma^1\cup\Sigma^2\cup\cdots$ 与 $\sigma=\sigma^1$，
便得到 $N$ 中端点落在 $\Sigma^0$ 内的路径的空间 $\scr{P}$ 上的一个函数 $I$ 与一个%
\trLineBundle{} $\delta$。这里有一点小麻烦：$\Sigma$ 是\trSingular{}的（\trCodimension{}
为二），而 $\sigma$ 只定义在它的\trSmooth{}部分上。为解决这一点，可以先对停留在
$\Sigma^0\cup\Sigma^1$ 内的路径 $\alpha$ 定义 $\delta$，然后再经一个涉及下一\trStratum{}
$\Sigma^2$ 的\trLink{}的计算，证明它可以延拓到整个 $\scr{P}$。细节略去。
与\trDeterminantLine{}\trLineBundle{}的联系经由\textbf{\trSpectralFlow}（spectral flow）构造而来。在 $\RR^n$
上取定一个\trInnerProduct，并取定一个函数 $\psi\in C^{\infty}(\RR,[0;1])$，满足：
当 $s\ll 0$ 时 $\psi(s)=0$，当 $s\gg 0$ 时 $\psi(s)=1$，而在一切使
$\psi(s)\in(0;1)$ 的 $s$ 处 $\psi'(s)>0$。对每个 $\alpha\in\scr{P}$，我们便得到一个%
\trFredholmOperator{} $d_{\alpha}=d/ds+\alpha(\psi(s))\in
\scr{F}(W^{1,2}(\RR,\RR^n),L^2(\RR,\RR^n))$。
\begin{lemma}\label{thm:spectral-flow-index}
$I(\alpha)=\mathrm{index}(d_{\alpha})$；而且 $\delta$ \trCanonical\trIsomorphic{}于%
\trFamily{} $\{d_{\alpha}\}$ 的\trDeterminantLine{}\trLineBundle。
\end{lemma}
\begin{proof}
取 $\alpha\in\scr{P}$ 使之与诸\trStratum{} $\Sigma^d$ \trTransverse，这特别地意味着
$\alpha^{-1}(\Sigma^d)=\varnothing$（对一切 $d\ge 2$）。用 $\nu_{\alpha(s)}$ 表示对称
矩阵 $\alpha(s)$ 的负\trEigenspace。先把 $d_{\alpha}$ 的\trDomain{}扩大：添上函数
$e^{-s\lambda}v$ 的\trLinear{}张成，其中 $\lambda$ 是 $\alpha(1)$ 的任一负\trEigenvalue，
而 $v$ 是相应的\trEigenvector。所得的算子
\begin{equation*}
W^{1,2}(\RR,\RR^n)\oplus\nu_{\alpha(1)}\longrightarrow L^2(\RR,\RR^n)
\end{equation*}
是\trSurjective{}的（因为可以从 $-\infty$ 积分到 $s$ 来求\trPreimage），而它的\trKernel{}通过在
靠近 $-\infty$ 的某点处取值可以与 $\nu_{\alpha(0)}$ 等同。这就给出
\begin{equation}\label{eq:kernel-index-dim}
\begin{gathered}
\dim\nu_{\alpha(0)}=\mathrm{index}(d_{\alpha})+\dim\nu_{\alpha(1)},\\
\lambda^{\mathrm{top}}(\nu_{\alpha(0)})\cong\det(d_{\alpha})\otimes
\lambda^{\mathrm{top}}(\nu_{\alpha(1)}),
\end{gathered}
\end{equation}
其中\trIsomorphism{}是\trCanonical{}的。在 $\alpha$ 穿越 $\Sigma^1$ 的有限多个时刻
$s_j$（$1\le j\le m$）的每一处，都有一个一维\trNullSpace{} $\sigma_{\alpha(s_j)}$，
并配有\trNondegenerate\trQuadraticForm{} $q_{\alpha}(s_j)$——它无非是 $\alpha'(s_j)$
的限制。容易看出存在\trIsomorphism
\begin{equation*}
\begin{cases}
\sigma_{\alpha(s_j)}\oplus\nu_{\alpha(s_j+\epsilon)}\cong\nu_{\alpha(s_j-\epsilon)}
  & \text{若 } q_{\alpha}(s_j)\text{ 为正，或}\\[2pt]
\nu_{\alpha(s_j+\epsilon)}\cong\sigma_{\alpha(s_j)}\oplus\nu_{\alpha(s_j-\epsilon)}
  & \text{若它为负，}
\end{cases}
\end{equation*}
其中 $\epsilon>0$ 充分小。在所有其余时刻，$\nu_{\alpha(s)}$ 随 $s$ \trSmooth{}变化；
把这些变化一路跟下来便得到
\begin{equation*}
\begin{gathered}
\dim\nu_{\alpha(0)}=\sum_j\mathrm{sign}(q_{\alpha}(s_j))+\dim\nu_{\alpha(1)},\\
\lambda^{\mathrm{top}}(\nu_{\alpha(0)})\cong\bigotimes_j
\sigma_{\alpha(s_j)}^{\mathrm{sign}(q_{\alpha}(s_j))}\otimes
\lambda^{\mathrm{top}}(\nu_{\alpha(1)}).
\end{gathered}
\end{equation*}
它与\eqref{eq:kernel-index-dim} 合在一起，给出所要的等式与\trIsomorphism。由于这一切
都是\trLinear{}代数，容易验证该\trIsomorphism{}与来自\eqref{eq:birth-isomorphism} 的
诸\trParallelTransport{}映射\trCompatible。
\end{proof}
现在我们把上面的论证重新表述成一种更适于无穷维推广的形式。首先，注意 $\delta_{\alpha}$
与 $d_{\alpha}$ 的定义可以平凡地推广到由任意闭区间 $C$ \trParametrization{}的路径 $\alpha$ 上。在一
点 $s\in\mathrm{int}(C)$（满足 $\alpha(s)\notin\Sigma$）处切开，把 $C$ 分解为
$C_1\cup C_2$。按定义，$I(\alpha)=I(\alpha|_{C_1})+I(\alpha|_{C_2})$，并且存在%
\trCanonical{}\trIsomorphism{} $\delta_{\alpha}\cong\delta_{\alpha|_{C_1}}\otimes
\delta_{\alpha|_{C_2}}$。另一方面，从\eqref{eq:kernel-index-dim} 可以看出
\begin{equation}\label{eq:cut-decomposition}
\begin{gathered}
\mathrm{index}(d_{\alpha})=\mathrm{index}(d_{\alpha|_{C_1}})+
\mathrm{index}(d_{\alpha|_{C_2}}),\\
\det(d_{\alpha})\cong\det(d_{\alpha|_{C_1}})\otimes\det(d_{\alpha|_{C_2}}).
\end{gathered}
\end{equation}
经一个不改变诸交点 $\alpha(s_j)$ 及该处相交符号的小\trDeformation，可以假设对每个
$j$ 都存在一个邻域 $[s_j-\epsilon;s_j+\epsilon]$，使得诸 $\alpha(s)$ 在其中同时\trDiagonalizable{}。把 $[0;1]$ 分解成如下区间：
\begin{equation*}
\begin{gathered}
C_0=[0;s_1-\epsilon],\\
C_1=[s_1-\epsilon;s_1+\epsilon],\\
C_2=[s_1+\epsilon;s_2-\epsilon],\\
\vdots\\
C_{2m}=[s_m+\epsilon;1].
\end{gathered}
\end{equation*}
多亏了同时对角化，诸算子 $d_{\alpha|_{C_{2j-1}}}$ 可以化归为纯量算子并作显式分析。
结论是：其中每一个要么\trInjective{}要么\trSurjective，带有一维的\trKernel{}或%
\trCokernel（视相交的符号而定），后者可以直接与 $\sigma_{\alpha(s_j)}$ 等同。对其余的
区间，$d_{\alpha|_{C_{2j}}}$ 有一个取定的\trDeformation{}通向一个\trInvertible{}算子：
只需在 $N\setminus\Sigma$ 内把 $\alpha|_{C_{2j}}$ \trRetract{}成取值为 $\alpha(s)$（对某个
$s\in C_{2j}$）的常值路径。为结束论证，应用\eqref{eq:cut-decomposition} 把各块重新
拼合起来即可。
